\begin{helper}[Completion of the two order-nine blocks]
Let $G$ be a finite loopless multigraph with individual edge identities,
let $\psi$ be a partial $k$-edge-coloring, and let $Q$ be balanced normal
globally maximum data for the tree-sequence
\[
 X=(y_0,a_0,y_1,a_1,y_2,a_2,y_3,a_3,y_4,
       b_0,y_5,b_1,y_6,c_0,y_7,c_1,y_8),
 \qquad a_0=e_X.
\]
Assume $X.\mathsf{length}=8$, $Q.\mathsf{height}=5$, and $h(G)=5$.
Write $\alpha=Q.\mathsf{alpha}$ and $\beta=Q.\mathsf{beta}$.
Let $d_0,d_1$ be colors, and let $m\in\{0,1,2,3,4\}$, with
\[
 d_0\in\overline\psi(y_0),\qquad
 d_1\in\overline\psi(y_m),\qquad
 d_0,d_1\notin\{\alpha,\beta\},\qquad d_1\ne d_0.
\]
Fix individual identities $x_1,x_2,\ell_B,g_0,g_A,g_B$ and vertices
$z_1,z_2$ satisfying all the following exact endpoint and color equations:
\[
\begin{array}{c|c|c}
\text{identity}&\operatorname{ends}&\psi\text{-color}\\ \hline
x_1&\{y_1,z_1\}&d_0\\
x_2&\{y_4,z_2\}&d_0\\
\ell_B&\{z_1,z_2\}&\beta\\
g_0&\{y_7,y_8\}&d_0\\
g_A&\{y_7,y_8\}&\alpha\\
g_B&\{y_7,y_8\}&\beta
\end{array}
\]
Assume either
\[
 (b_0,y_5,b_1,y_6)=(x_1,z_1,x_2,z_2)
 \quad\text{or}\quad
 (b_0,y_5,b_1,y_6)=(x_2,z_2,x_1,z_1),
\]
and assume $\psi(c_0)=\psi(c_1)=d_1$.
Then there exist witnesses that $Q$ is balanced,
$X.\mathsf{length}=8$, and $Q.\mathsf{height}=5$, and an injective map
$\gamma:\{0,1\}\to[k]$, such that for every $i,j\in\{0,1\}$:
\begin{enumerate}
\item the exact identity $X.\mathsf{edge}(4+2i+j)$ has endpoints
      $\{y_q,y_{5+2i+j}\}$ for some $q\in\{0,1,2,3,4\}$;
\item $\psi(X.\mathsf{edge}(4+2i+j))=\gamma(i)$;
\item $\gamma(i)$ is missing at some $y_q$ with $0\leq q<5$;
\item every color on $a_1,a_2,a_3$ differs from both
      $\gamma(0)$ and $\gamma(1)$.
\end{enumerate}
The conclusion concerns the same coloring, exact ordered tree, and normal
data $\psi,X,Q$, not a replacement presentation.
\end{helper}
\begin{proof}
The global field of \noderef{NormalGlobalMaximumTriple} supplies
properness in the sense of \noderef{IsProperEdgeColoring} and the
unique-uncolored-edge condition \noderef{IsOneEdgeUncolored}.
Vertex and edge injectivity in \noderef{TreeSequence} make the nine
displayed vertices distinct and the eight displayed identities distinct.
Either first-block orientation gives $\{z_1,z_2\}=\{y_5,y_6\}$,
with $z_1\ne z_2$.  The exact normal-cycle relations of
\noderef{ScheideTrunkRotationCycle} give
\[
 \psi(a_1)=\psi(a_3)=\alpha,\quad \psi(a_2)=\beta,
 \quad\operatorname{ends}(q)=\{y_4,y_0\},\quad\psi(q)=\beta,
 \qquad q=Q.\mathsf{closingEdge}.
\]
Also $\alpha$ is missing at $y_0$ and $\beta$ at $y_1$.

We first prove that one of $c_0,c_1$ joins the trunk to its newly
displayed vertex.  Suppose neither does.  The attachment equations put
the earlier endpoint of $c_0$ among $y_0,\ldots,y_6$ and that of $c_1$
among $y_0,\ldots,y_7$.  By the supposition neither earlier endpoint is
on the trunk.  The latter endpoint cannot be $y_7$: both distinct
identities would then have color $d_1$ at $y_7$, contrary to properness.
Thus both earlier endpoints are among $y_5,y_6$, and they must be
different, again by properness.  There is consequently a labeling
$\{a,b\}=\{y_7,y_8\}$ and exact identities
$\{r_5,r_6\}=\{c_0,c_1\}$ with
\[
 \operatorname{ends}(r_5)=\{z_1,a\},\quad
 \operatorname{ends}(r_6)=\{z_2,b\},\quad
 \psi(r_5)=\psi(r_6)=d_1.
\]
This covers both orientations of each pair.

In \noderef{TwoColorIncidenceGraph} for $d_1,\beta$, the four
identities $\ell_B,r_6,g_B,r_5$ give the cycle
\[
 z_1,\ell_B,z_2,r_6,b,g_B,a,r_5,z_1.
\]
Its four vertices are distinct.  The two $d_1$ identities are distinct
by tree-edge injectivity, and each differs from each $\beta$ identity
by the color inequality.  The two $\beta$ identities have disjoint
endpoint pairs and are therefore distinct as well.  Every vertex-node
already has two different neighbors in this cycle; properness and
\noderef{TwoColorIncidenceDegreeAtMostTwo} allow no further neighbor.
Each edge-node has exactly its two loopless endpoints, by
\noderef{Incident}.  Hence no adjacency leaves the displayed cycle.
Its nodes are connected, so they are exactly the
\noderef{TwoColorComponent} rooted at $z_1$.

Swap this full component using \noderef{TwoColorComponentSwap} and call
the new coloring $\chi$.  It is proper by
\noderef{TwoColorComponentSwapProper}.  The cycle avoids every trunk
vertex, so every trunk identity, including the uncolored identity and
$q$, is fixed.  All other identities of the displayed tree except
$c_0,c_1$ are also fixed: $x_1,x_2$ have the excluded color $d_0$,
and the other displayed identities lie on the trunk.  The last pair now
has color $\beta$.  By \noderef{TwoColorComponentSwapMissingBehavior},
$d_0$ and $\alpha$ are still missing at $y_0$, and $\beta$ is still
missing at $y_1$, which is outside the swapped component.
These witnesses verify \noderef{TashkinovTree} for the same sequence:
$a_1,a_3$ use $\alpha$ at $y_0$, $a_2$ uses $\beta$ at $y_1$,
the first pair uses $d_0$ at $y_0$, and the last pair uses $\beta$ at
$y_1$.  The swap preserves the unique uncolored identity and all other
edges remain colored.  The coloring-independent triple hypotheses and
the vertex cardinality are unchanged.  Thus the same tree with $\chi$
satisfies \noderef{GlobalMaximumTriple}.

In the $d_0,\beta$ incidence graph of $\chi$, the exact identities give
the connected sequence
\[
 y_1,x_1,z_1,r_5,a,g_0,b,r_6,z_2,x_2,y_4,q,y_0.
\]
Here $x_1,x_2,g_0$ retain color $d_0$, $r_5,r_6$ now have color
$\beta$, and $q$ retains color $\beta$.  Apply
\noderef{ScheideKempePathHeightBound} with $d_0$ missing at $y_0$
and $\beta$ missing at $y_1$.  It gives a simple incidence path $P$
between these vertices, whose support is their full two-color component,
with $4\mid L$ and $L/2+1\leq h(G)=5$, where $L$ is its incidence
length.  Every node of the connected sequence is therefore in its
support.  In particular, its support contains seven distinct vertex-nodes
$y_0,y_1,y_4,z_1,z_2,a,b$ and four distinct edge-nodes
$x_1,x_2,r_5,r_6$.  The first two identities are distinct by the given
orientation and tree-edge injectivity; the last two are distinct for the
same reason; and identities from different pairs have different original
colors $d_0,d_1$.  Vertex-nodes and edge-nodes are disjoint kinds.
Consequently $L+1=|\operatorname{support}(P)|\geq11$.  But the height
bound and divisibility imply $L\leq8$, a contradiction.  This proves
the claimed trunk crossing.

Next consider a presentation in which $d_1$ is missing at its first
vertex.  The known crossing and
\noderef{ScheideHeightFiveGammaEdges} supply two exact $d_1$ crossings
from trunk positions one and four to outside vertices $u_1,u_2$, and
one internal $d_1$ identity on trunk positions two and three.
The known crossing cannot attach at position zero, by missingness, or
at positions two or three, since properness there would identify it
with the internal identity and force its outside endpoint onto the
trunk.  Properness at position one or four instead identifies it with
the corresponding supplied crossing.  Thus at least one of $u_1,u_2$
is $y_7$ or $y_8$.
By \noderef{ScheideHeightFiveOutsideLinks}, an identity of the first
normal color joins $u_1,u_2$.  At their known last-pair endpoint,
properness identifies that identity with the specified last-pair
identity of the same color.  Equality of endpoint pairs then gives
$\{u_1,u_2\}=\{y_7,y_8\}$.  At each last vertex, properness
identifies its incident displayed $d_1$ tree identity with the supplied
$d_1$ crossing.  Both $c_0,c_1$ therefore attach to the trunk.

If $m=0$, apply this argument directly.  If $m>0$, use
\noderef{ScheideTrunkRotationColoring},
\noderef{ScheideTrunkRotationTree}, and
\noderef{ScheideTrunkRotationNormal} to rotate the trunk so that $y_m$
is first.  Their exact formulas preserve all post-trunk vertices and
identities and their colors, rotate the trunk vertices bijectively, and
preserve height five and balancedness.  Only identities on the trunk
cycle change, and every new colored value there is a trunk color.
Since $d_1$ differs from both trunk colors, a hypothetical new
$d_1$ incidence at the moved missing vertex would either be on that
cycle, contradicting this color exclusion, or be unchanged, contradicting
the old missingness.  Thus $d_1$ is missing at the new first vertex.
All three exact last-pair identities lie outside the trunk cycle and are
fixed.  The temporary trunk colors are $\alpha,\beta$ in either order,
so the first-normal-color last-pair identity is $g_A$ or $g_B$ as
appropriate.  The known crossing transports using the inverse trunk
permutation.  Applying the preceding localization and then this
permutation back proves both required crossings in the original
presentation, without changing $\psi,X,Q$.

Finally set $\gamma(0)=d_0$ and $\gamma(1)=d_1$.  The inequality
$d_1\ne d_0$ proves injectivity.  The given first-block orientation
proves its two endpoint and color assertions; the localization proves
the second-block endpoint assertions, and its colors were assumed.
The missing witnesses are positions zero and $m$.  The normal identities
$a_1,a_3$ have color $\alpha$ and $a_2$ has color $\beta$, so all
three avoid both block colors by the assumed exclusions.  Retain the
given balancedness, length, and height witnesses.  This is the full
asserted conclusion.
\end{proof}
